% !TEX program = lualatex
\documentclass[12pt, a4paper]{article}
\usepackage{master}
\usepackage{longtable}

% 마진 축소
\newgeometry{
  top=18mm, bottom=16mm, left=15mm, right=15mm,
  headsep=5mm, footskip=10mm
}

\fancyhead[L]{\footnotesize\color{midgray}오르가논 강독}
\fancyhead[R]{\footnotesize\color{midgray}42주 스케줄}

\begin{document}
\thispagestyle{firstpage}

\section{오르가논 42주 강독 스케줄}

아리스토텔레스 오르가논 여섯 저작과 포르피리오스 이사고게의 42주 강독 계획입니다. 주당 분량은 베커판 기준 약 4--7쪽이며, 복습 주간 없이 전체를 강독에 배정합니다.

\vspace{6pt}

{%
\footnotesize
\setlength{\extrarowheight}{1pt}%
\renewcommand{\arraystretch}{1.05}%
\setlength{\tabcolsep}{3pt}%
\rowcolors{2}{altrowbg}{white}%
\begin{longtable}{|>{\centering\arraybackslash}m{6mm}|>{\centering\arraybackslash}m{30mm}|>{\centering\arraybackslash}m{24mm}|>{\centering\arraybackslash}m{102mm}|}
\hline
\rowcolor{tableheadblue}
{\color{h1navy}\bfseries 주} &
{\color{h1navy}\bfseries 저작} &
{\color{h1navy}\bfseries 범위} &
{\color{h1navy}\bfseries 주제} \\ \hline
\endfirsthead
\hline
\rowcolor{tableheadblue}
{\color{h1navy}\bfseries 주} &
{\color{h1navy}\bfseries 저작} &
{\color{h1navy}\bfseries 범위} &
{\color{h1navy}\bfseries 주제} \\ \hline
\endhead
\hline
\endfoot
%
% ── 오리엔테이션 ──
\rowcolor{reviewbg}
1 & 오리엔테이션 & — & 오르가논의 이름, 여섯 저작의 배열, 수용사 \\ \hline
%
% ── 범주론 (2–5주차) ──
2 & 범주론 & 1a1--4a21 & 동음이의·동의·파생어, 사중 분류, 10범주, 실체(전반) \\ \hline
3 & 범주론 & 4a22--6b35 & 반대자 수용 논변 완결, 양 범주 \\ \hline
4 & 범주론 & 6b36--11a38 & 관계 범주(두 정의), 질 범주(네 종류) \\ \hline
5 & 범주론 & 11a38--15b32 & 나머지 범주, 후범주론(대립·선후·동시·변화) \\ \hline
%
% ── 이사고게 (6–8주차) ──
6 & 이사고게 & I--III & 서론, 류, 종 \\ \hline
7 & 이사고게 & IV--V & 종차, 고유성 \\ \hline
8 & 이사고게 & VI--VII & 부수성, 다섯 술어가능자의 비교 \\ \hline
%
% ── 명제론 (9–11주차) ──
9 & 명제론 & 16a1--17b16 & 이름·동사·문장, 긍정과 부정, 단칭·보편 명제 \\ \hline
10 & 명제론 & 17b16--19a22 & 대당 관계, 양상 명제 \\ \hline
11 & 명제론 & 19a23--21b33 & 해전 논변(9장), 복합 명제, 요약 \\ \hline
%
% ── 분석론 전서 (12–22주차) ──
12 & 분석론 전서 & 24a10--26b33 & 삼단논법 정의, 제1격 \\ \hline
13 & 분석론 전서 & 26b34--29a18 & 제2격, 제3격 \\ \hline
14 & 분석론 전서 & 29a19--31b12 & 격의 비교, 부정 삼단논법 \\ \hline
15 & 분석론 전서 & 31b12--33b17 & 양상 삼단논법 도입, 필연 전제 \\ \hline
16 & 분석론 전서 & 33b17--36a17 & 가능 전제, 혼합 양상 \\ \hline
17 & 분석론 전서 & 36a17--39a3 & 혼합 양상(계속), 환원법 \\ \hline
18 & 분석론 전서 & 39a4--42a31 & 용어 선택, 전환 규칙, 연쇄 삼단논법 \\ \hline
19 & 분석론 전서 & 42a32--45b20 & 추론의 발견법 \\ \hline
20 & 분석론 전서 & 45b21--50a4 & 순환 논증, 귀류법 \\ \hline
21 & 분석론 전서 & 50a5--62b38 & 오류 추론, 반론 \\ \hline
%
% ── 페이지 분리: 1–21주차 / 22–42주차 ──
\pagebreak
%
22 & 분석론 전서 & 62b39--70b38 & II권: 역전, 귀납, 예증, 반론, 개연적 추론 \\ \hline
%
% ── 분석론 후서 (23–29주차) ──
23 & 분석론 후서 & 71a1--73b24 & 학문적 인식, 논증의 정의, 자체 술어 \\ \hline
24 & 분석론 후서 & 73b25--76b22 & 논증의 원리, 원리의 종류, 공리 \\ \hline
25 & 분석론 후서 & 76b23--79b22 & 학문의 통일성, 학문 간 관계 \\ \hline
26 & 분석론 후서 & 79b23--84a6 & 무한 소급 불가, 논증의 유한성 \\ \hline
27 & 분석론 후서 & 84a7--87b22 & 필연성, 보편성, 오류 \\ \hline
28 & 분석론 후서 & 87b23--93a9 & II권: 정의와 논증의 관계 \\ \hline
29 & 분석론 후서 & 93a10--100b17 & 정의의 종류, 분류법, 누스(II.19) \\ \hline
%
% ── 토피카 (30–39주차) ──
30 & 토피카 & 100a18--104b17 & I권: 변증술의 목적, 통념, 명제, 문제 \\ \hline
31 & 토피카 & 104b18--114a25 & II권: 부수성에 관한 토포이 \\ \hline
32 & 토피카 & 114a26--122b39 & III권: 비교의 토포이 \\ \hline
33 & 토피카 & 120a15--128b11 & IV권: 류에 관한 토포이 \\ \hline
34 & 토피카 & 128b12--137a7 & V권: 고유성에 관한 토포이 \\ \hline
35 & 토피카 & 137a8--145b20 & VI권: 정의에 관한 토포이(전반) \\ \hline
36 & 토피카 & 145b21--155a2 & VI권(후반)--VII권: 정의(계속), 동일성 \\ \hline
37 & 토피카 & 155a3--160b16 & VIII권: 변증술적 논쟁의 규칙(전반) \\ \hline
38 & 토피카 & 160b17--164b19 & VIII권(후반): 변증술적 논쟁의 전략 \\ \hline
39 & 토피카 & — & 토피카 종합: 변증술의 체계적 의의 \\ \hline
%
% ── 소피스트적 논박 (40–42주차) ──
40 & 소피스트적 논박 & 164a20--172b8 & 소피스트적 논박의 목적, 언어 의존 오류 \\ \hline
41 & 소피스트적 논박 & 172b9--180b7 & 언어 비의존 오류, 오류 해소법 \\ \hline
\rowcolor{reviewbg}
42 & 소피스트적 논박 & 180b8--184b8 & 오류 해소(계속), 아리스토텔레스의 에필로그, 총정리 \\ \hline
\end{longtable}
}

\end{document}
